%! Solving Rational Equations — Unit 4, Lesson 4.6 — Homework (Answer Key)
\documentclass[10pt]{article}
\usepackage{algebra2-article}
\usepackage{algebra2-key}   % pulls in -boxes; do NOT also load -boxes
\newcommand{\TallMath}[1]{$\displaystyle #1\rule[-1.4em]{0pt}{3.2em}$}
% Standard coordinate grid + axes: {xmin}{xmax}{ymin}{ymax}
\newcommand{\lgrid}[4]{%
  \draw[step=1, linegray!40, very thin] (#1,#3) grid (#2,#4);
  \draw[->, semithick, charcoal] (#1,0)--(#2,0) node[below right, font=\scriptsize]{$x$};
  \draw[->, semithick, charcoal] (0,#3)--(0,#4) node[left, font=\scriptsize]{$y$};}

\begin{document}

\pageheader{Unit 4, Lesson 4.6}{Homework --- Answer Key}
\namedateperiod

\vspace{0.10in}

\begin{notesbox}{Practice}
Every problem starts the same way: \textbf{factor the denominators and write the restriction list.}
A number is not an answer until you have checked it against that list.

\begin{enumerate}[leftmargin=1.9em, itemsep=10pt]
  \item \textbf{The list first.} Do not solve --- just give the restrictions and the LCD.
    \begin{enumerate}[label=(\alph*), leftmargin=2.1em, itemsep=4pt]
      \item $\dfrac{7}{x}-\dfrac{2}{x+5}=1$ \hfill restrictions \ans{$0,\,-5$} \; LCD \ans{$x(x+5)$}
      \item $\dfrac{x}{x-7}=\dfrac{3}{x+1}$ \hfill restrictions \ans{$7,\,-1$} \;
            LCD \ans{$(x-7)(x+1)$}
      \item $\dfrac{4}{x^2-25}=\dfrac{1}{x-5}$ \hfill restrictions \ans{$5,\,-5$} \;
            LCD \ans{$(x-5)(x+5)$}
      \item $\dfrac{2}{x^2+2x}+\dfrac{1}{x+2}=\dfrac{3}{x}$ \hfill restrictions \ans{$0,\,-2$} \;
            LCD \ans{$x(x+2)$}
    \end{enumerate}

  \item \textbf{Solve.} Show the cleared equation, the candidates, and the check. If an equation has no
        solution, say so.
    \begin{enumerate}[label=(\alph*), leftmargin=2.1em, itemsep=7pt]
      \item $\dfrac{5}{x}+\dfrac{3}{2}=\dfrac{11}{x}$ \hfill candidates \ans{$4$} \;
            extraneous \ans{none} \; solution \ans{$\{4\}$}
      \item $\dfrac{x}{x+3}=\dfrac{4}{x}$ \hfill candidates \ans{$6,\,-2$} \;
            extraneous \ans{none} \; solution \ans{$\{6,-2\}$}
      \item $\dfrac{x}{4}+\dfrac{1}{x}=\dfrac{5}{4}$ \hfill candidates \ans{$1,\,4$} \;
            extraneous \ans{none} \; solution \ans{$\{1,4\}$}
      \item $\dfrac{x}{x-6}=\dfrac{6}{x-6}+2$ \hfill candidates \ans{$6$} \;
            extraneous \ans{$6$} \; solution \ans{none}
      \item $\dfrac{x}{x-1}+\dfrac{1}{x+1}=\dfrac{2}{x^2-1}$ \hfill candidates \ans{$-3,\,1$} \;
            extraneous \ans{$1$} \; solution \ans{$\{-3\}$}
      \item $\dfrac{4}{x^2-16}+\dfrac{1}{x+4}=\dfrac{1}{x-4}$ \hfill candidates \ans{none} \;
            extraneous \ans{---} \; solution \ans{none}
    \end{enumerate}
    \vspace{3pt}
    Parts (d) and (f) both come out to ``no solution,'' but for two different reasons. Explain the
    difference. \ansline{In (d) the cleared equation had a perfectly good candidate, $x=6$, and it was
    thrown out only because $6$ is an excluded value --- an extraneous solution.}
    \ansline{In (f) the cleared equation reduces to $0=4$, which is never true, so there was no
    candidate to throw out in the first place. No extraneous solution is involved.}

  \item \textbf{Check, do not solve.} A candidate is given. Decide, and give the reason.
    \begin{enumerate}[label=(\alph*), leftmargin=2.1em, itemsep=5pt]
      \item $\dfrac{x+2}{x-4}=\dfrac{6}{x-4}$, \; candidate $x=4$: \ans{no}, because
            \ans{$x-4=0$ there; both sides are undefined --- extraneous}
      \item $\dfrac{3}{x}-\dfrac{1}{x-2}=\dfrac{1}{x}$, \; candidate $x=4$: \ans{yes}, because
            \ans{$\frac34-\frac12=\frac14$, and $4$ is not excluded}
      \item $\dfrac{5}{x+1}=\dfrac{x+3}{x+1}$, \; candidate $x=2$: \ans{yes}, because
            \ans{both sides equal $\frac53$, and $2$ is not excluded}
    \end{enumerate}

  \item \textbf{See it on a graph.} Below is the graph of the single fraction you get by moving
        everything in problem \textbf{2(e)} to one side and combining it.
    \begin{center}
    \begin{tikzpicture}[scale=0.36]
      \lgrid{-8}{7}{-5}{7}
      \draw[navy, dashed, semithick] (-1,-5)--(-1,7);
      \draw[navy, dashed, semithick] (-8,1)--(7,1);
      \node[navy, font=\scriptsize, above left] at (-1,6.2) {$x=-1$};
      \node[navy, font=\scriptsize, below right] at (4.6,1) {$y=1$};
      \begin{scope}
        \clip (-8,-5) rectangle (7,7);
        \draw[very thick, forest, domain=-8:-1.35, samples=140, smooth] plot (\x,{1+2/((\x)+1)});
        \draw[very thick, forest, domain=-0.6:7, samples=140, smooth] plot (\x,{1+2/((\x)+1)});
      \end{scope}
      \draw[forest, very thick, fill=white] (1,2) circle (4.2pt);
      \fill[charcoal] (-3,0) circle (3.2pt) node[below left, font=\scriptsize]{$(-3,0)$};
    \end{tikzpicture}
    \end{center}
    \vspace{-4pt}
    \begin{enumerate}[label=(\alph*), leftmargin=2.1em, itemsep=5pt]
      \item Combine and simplify: \;
            $\dfrac{x}{x-1}+\dfrac{1}{x+1}-\dfrac{2}{x^2-1}
            =\dfrac{{\color{keyred}\mathbf{(x+3)(x-1)}}}{{\color{keyred}\mathbf{(x-1)(x+1)}}}
            ={\color{keyred}\mathbf{\dfrac{x+3}{x+1}}}$, \;
            $x\neq{\color{keyred}\mathbf{1}}$
      \item The solution of the equation is at the graph's \ans{$x$-intercept}, namely $x=$ \ans{$-3$}.
      \item The extraneous candidate is at the graph's \ans{hole}, namely $x=$ \ans{$1$}.
      \item There is a third interesting $x$-value here: $x=-1$, where the graph has a
            \ans{vertical asymptote}. Was $-1$ ever a candidate? \ans{no} \; So: is every excluded value
            an extraneous solution? \ans{no}
    \end{enumerate}

  \item \textbf{Error analysis.} A student solved $\dfrac{10}{x}-\dfrac{4}{x-1}=1$ and wrote:
    \begin{center}
    \fbox{\parbox{0.82\linewidth}{\centering \vspace{2pt}
    ``LCD is $x(x-1)$. Multiply the fractions by it: \\
    $10(x-1)-4x=1$, \; so $6x-10=1$, \; so $x=\frac{11}{6}$. \\
    It's not $0$ or $1$, so it checks out.''\vspace{2pt}}}
    \end{center}
    \begin{enumerate}[label=(\alph*), leftmargin=2.1em, itemsep=5pt]
      \item The LCD is correct and the check at the end is correct. The mistake is in the
            \ans{multiply-by-the-LCD} step. Describe it in a few words. \ans{they multiplied only the
            two fractions, not both entire sides}
      \item Which term did the student fail to multiply? \ans{the $1$} \; What should the right side
            have become? \ans{$x(x-1)$}
      \item Redo it correctly. Cleared equation: \ans{$6x-10=x^2-x$}. \; Candidates \ans{$2,\;5$}.
      \item Extraneous \ans{none}. \; Solution set \ans{$\{2,\,5\}$}.
    \end{enumerate}

  \item \textbf{Paddling against the current.} A kayaker paddles at a steady $4$ miles per hour in
        still water. On a river with a current of $c$ miles per hour, she goes $6$ miles
        \emph{upstream} and then the same $6$ miles back \emph{downstream}. The whole trip takes
        $4$ hours.
    \begin{enumerate}[label=(\alph*), leftmargin=2.1em, itemsep=5pt]
      \item Going upstream her speed over the ground is \ans{$4-c$} mph, so that leg takes
            \ans{$\frac{6}{4-c}$} hours. Going downstream her speed is \ans{$4+c$} mph, so that leg takes
            \ans{$\frac{6}{4+c}$} hours.
      \item Write the equation for the whole trip: \; \ans{$\dfrac{6}{4-c}+\dfrac{6}{4+c}=4$}
      \item Restrictions: $c\neq$ \ans{$4,\;-4$}. \; In plain words, what would a current of exactly
            $4$ mph do to the upstream leg, and how does the graph of the upstream time show it?
            \ansline{the current would exactly cancel her paddling, so she never moves upstream and the
            time is undefined --- the graph of $\frac{6}{4-c}$ has a vertical asymptote at $c=4$}
      \item Solve. Candidates \ans{$2,\;-2$}. Are either of them extraneous? \ans{no}
      \item You must still reject one of them. Which, and \textbf{why}? Is your reason the same kind of
            reason as ``extraneous''? \ansline{Reject $c=-2$: a current speed cannot be negative. This
            is a \emph{context} rejection, not an algebraic one.}
            \ansline{$c=-2$ really does satisfy the equation, unlike an extraneous root, which
            satisfies only the cleared equation and makes the original undefined.}
      \item The current is \ans{$2$} mph. Check it: upstream takes \ans{$3$} hours,
            downstream takes \ans{$1$} hours, total \ans{$4$} hours.
    \end{enumerate}
\end{enumerate}
\end{notesbox}

\begin{extensionbox}
\textbf{Part 1 --- Does every excluded value get a turn?} Problem 4(d) suggests an answer. Write a
rational equation whose restriction list includes $x=2$, but for which $x=2$ never appears as a
candidate at all.
\begin{center}
\ans{e.g.\ $\dfrac{1}{x-2}=1$, whose only candidate is $x=3$ \; (answers vary)}
\end{center}
\vspace{-2pt}
Then finish the sentence: every extraneous solution is an excluded value, but \ans{not every excluded
value turns up as a candidate --- most of them never appear in the cleared equation at all}

\vspace{5pt}
\textbf{Part 2 --- One number decides everything.} Consider the family
\[ \frac{x}{x-2}=\frac{k}{x-2}, \]
where $k$ is a number you get to choose.
\begin{enumerate}[label=(\alph*), leftmargin=2.1em, itemsep=5pt]
  \item Clear the denominator. The cleared equation is \ans{$x=k$}, so the only candidate is
        $x=$ \ans{$k$}.
  \item For exactly one value of $k$, the equation has \emph{no} solution. That value is
        $k=$ \ans{$2$}. Why that one? \ans{then the only candidate is $2$, the excluded value}
  \item For every other $k$, the solution is $x=$ \ans{$k$}. Test your rule on $k=7$ and on
        $k=-5$. \ansline{$k=7$ gives $x=7$, and both sides equal $\frac75$; $k=-5$ gives $x=-5$, and
        both sides equal $\frac57$ --- neither is the excluded value $2$, so both are genuine}
\end{enumerate}
\end{extensionbox}

\begin{spiralbox}
\textbf{Coming next --- Lesson 4.7: Direct, Inverse \& Joint Variation.} The unit closes where it
started: with $y=\frac{k}{x}$, the rational parent, wearing its working clothes. You will read a table
and decide whether two quantities are \emph{directly} proportional (one doubles, so does the other) or
\emph{inversely} proportional (one doubles, the other halves), find the constant $k$, write the
equation, and say what it means. Today's habit carries straight over: when the answer has to describe
something real, some solutions get rejected --- and you now know the difference between rejecting one
because the algebra forbids it and rejecting one because the situation does.
\end{spiralbox}

\begin{teachernote}
\textbf{Problem 2 is the spine, and 2(d) versus 2(f) is the part worth grading closely.} Both answers
are ``no solution,'' and students will treat them as the same event. They are not. In (d) the cleared
equation produced $x=6$ and the check removed it --- an extraneous solution. In (f) the cleared
equation collapses to $0=4$; there was never a candidate, so nothing was extraneous. If a student
writes ``$x=6$ is extraneous'' for (f), they have pattern-matched rather than solved.

\vspace{3pt}
\textbf{Problem 4 is the graphical-verification item} (A2.EI.4b/c), and 4(d) is the conceptual payoff of
the whole assignment. The restriction list for 2(e) is $\{1,-1\}$; the candidates were $\{-3,1\}$. So
$-1$ is excluded and yet never extraneous --- it is the \emph{wall}, and the equation simply never
proposed it. The implication is one direction only: \textbf{every extraneous solution is an excluded
value, but not every excluded value is extraneous.} Students who believe the converse will start
throwing away good answers.

\vspace{3pt}
\textbf{Problem 5 targets the other big error of the day} --- not the missing check, but the incomplete
multiplication. The student's LCD is right, their restriction list is right, and their final check is
right; they multiplied the \emph{fractions} instead of both \emph{sides}. Correct work:
$10(x-1)-4x=x(x-1)$, so $x^2-7x+10=0$ and $x=2$ or $x=5$, both valid. It is worth writing the
$\cdot\,x(x-1)$ under every single term when demonstrating.

\vspace{3pt}
\textbf{Problem 6 is the modeling item} (A2.EI.4a/c). Students should build $\frac{6}{4-c}+\frac{6}{4+c}
=4$ themselves from $\text{time}=\frac{\text{distance}}{\text{rate}}$; do not hand them the equation.
Clearing gives $48=64-4c^2$, so $c=\pm2$ and \emph{neither is extraneous}. The rejection of $c=-2$ is a
context rejection --- part (e) is asking for exactly that distinction, and it is the part to look at
first when grading. Part (c) is a deliberate callback to 4.5: $c=4$ is a vertical asymptote of the
upstream-time function, and it means the kayaker is paddling exactly as fast as the river pushes back.

\vspace{3pt}
\textbf{Extension Part 2} generalizes the Hook. The candidate is always $x=k$, so the equation fails
exactly when $k$ lands on the excluded value $2$. Students who see this have understood that an
extraneous solution is a coincidence of position, not a property of the algebra.
\end{teachernote}

\end{document}
